\chapter{Metodología comparativa del estudio}
\label{ch:metodologia-comparativa}

\section{Introducción del capítulo}
Una comparación entre algoritmos solo es válida si todos enfrentan el mismo problema bajo
condiciones equivalentes. En estudios de sintonización PID esto significa utilizar la misma
planta, la misma función objetivo, los mismos límites para las ganancias y escenarios de prueba
igualmente exigentes. De lo contrario, las diferencias observadas podrían provenir del diseño
experimental y no del algoritmo evaluado.

\section{Objetivo del capítulo}
El objetivo de este capítulo es establecer un marco de comparación justo entre la sintonización
clásica y los métodos metaheurísticos desarrollados previamente. Se busca definir cómo se
construyen los contendientes, qué métricas se observan, qué escenarios se consideran y qué
criterios permiten interpretar los resultados sin sesgos metodológicos.

\section{Diseño experimental}
La lógica comparativa propuesta toma como línea base al controlador ajustado mediante
Ziegler--Nichols y enfrenta esa referencia con controladores PID sintonizados por PSO, GWO,
Eagle Strategy y ABC. En todos los casos, la planta modelada es la misma microred
hidro-solar de Limbani y la calidad del ajuste se mide a partir de indicadores transitorios y de
índices integrales de error. Esto garantiza que la comparación recaiga en la capacidad de
búsqueda de cada método y no en diferencias de modelado.

\section{Plantas o sistemas de prueba}
La planta principal es una microred híbrida con generación hidráulica de flujo cruzado y apoyo
fotovoltaico. Sobre esa base se consideran varias situaciones operativas: régimen nominal,
escalón positivo de carga, intermitencia solar y pruebas de estrés con restricciones adicionales.
Esta selección permite observar tanto el comportamiento del controlador cerca del punto de
operación como su capacidad para responder ante eventos críticos.

\section{Condiciones homogéneas de evaluación}
Para mantener equidad experimental, todos los algoritmos trabajan sobre el mismo rango de
ganancias, con idéntica función objetivo basada en ITAE y con igual horizonte de simulación.
Cuando la estructura del algoritmo lo permite, se procuran presupuestos de búsqueda
comparables en número de evaluaciones de la planta. Esta decisión evita favorecer a un método
solo por haberle concedido más intentos de exploración o un problema numérico diferente.

\section{Métricas de comparación}
El análisis final combina criterios cuantitativos y cualitativos. Entre los primeros destacan
sobreimpulso, tiempo de asentamiento, error estacionario, ITAE y, cuando corresponde, costo
computacional. Entre los segundos se consideran la suavidad de la actuación, la robustez frente
a perturbaciones distintas y la coherencia física de la solución encontrada. Esta lectura mixta es
necesaria porque una sintonización óptima en papel no siempre es la más conveniente para
proteger el equipo o facilitar una implementación real.

\section{Cierre del capítulo}
Con el marco comparativo definido, el siguiente paso consiste en describir con mayor detalle los
casos de estudio y los escenarios de operación que servirán para poner a prueba las distintas
estrategias de sintonización.
